% Part: proof-theory
% Chapter: normalization
% Section: introduction

\documentclass[../../../include/open-logic-section]{subfiles}

\begin{document}

\olfileid{pt}{nor}{int}

\olsection{جان پچھان}

جے تسی شرطیہ~$!A \lif !B$ !!{prove}d کر چکے او، تاں ہو سکدا اے کہ !!a{proof}~$!B$ دا دین ویلے~\Elim{\lif} لا کے ایہنوں ورت سکو۔ ایس لئی!!a{proof}~$\delta_1$،~$!A$ دا، وی چاہیدا ہووے گا۔ تہاڈا !!{proof}~$\delta_1$،~$!A \lif !B$ دا، غالباً~\Intro{\lif} اُتے مُکدا اے، جیہڑا !!a{proof}~$\delta_2$،~$!B$ دا کھُلّے مفروضے~$!A$ توں،~$!A \lif !B$ دے ثبوت وچ بدلدا اے۔ جے ایویں اے، تاں تہاڈے آخری !!{proof} دی شکل ایہ اے:
\[
\AxiomC{$\Discharge{!A}{x}$}
\RightLabel{$\delta_2$}
\DeduceC{$!B$}
\DischargeRule{\Intro{\lif}}{x}
\UnaryInfC{$!A \lif !B$}
\AxiomC{}
\RightLabel{$\delta_2$}
\DeduceC{$!A$}
\RightLabel{\Elim{\lif}}
\BinaryInfC{$!B$}
\DisplayProof
\]
بھاویں تہاڈا طریقہ ایس لحاظ نال چنگا اے کہ تسی پہلاں توں موجود شے مُڑ ورتی (!!a{proof}~$!A \lif !B$ دا)، تہاڈا !!{proof} اوہنا چنگا نہیں۔ $!A \lif !B$ نوں شامل کرنا تے جھٹ ہٹا دینا اک چکر لگدا اے۔~$!B$ نوں !!{prove} کرن دا ہور سدھا طریقہ ہونا چاہیدا اے۔

اصل وچ ایہو جیہا طریقہ ہمیشہ ہوندا اے۔ اُتلے ورگے «چکر»، جتھے !!a{introduction} قاعدے دے جھٹ پچھوں !!a{elimination} قاعدہ آؤندا اے، قدرتی اخذ دے !!{proof}s وچوں ہمیشہ ہٹائے جا سکدے نیں۔ ایس نوں \emph{نارمل روپ وچ لیاؤنا} آکھدے آں تے حاصل اک \emph{نارمل} !!{proof} (یا !!a{proof} \emph{نارمل روپ وچ}) اے۔

ایس نتیجے دا ثبوت، سیکوئنٹ حساب دے قطع ہٹاؤن والے قضیے وانگ، ذرا گنجل دار اے تے کئی صورتاں ویکھنیاں پیندیاں نیں۔ کلاسیکی \Log{N1c} لئی ثابت کرنا وجدانی \Log{N1i} نالوں اوکھا اے۔ جیویں قطع ہٹاؤنا وکھاؤندا اے کہ سیکوئنٹ حساب وچ ہور کجھ !!{prove} نہیں کر سکدے، \Cut{} دا قاعدہ ورت کے، اویویں نارمل روپ وچ لیاؤنا وکھاؤندا اے کہ چکر پا کے اوس توں ودھ نہیں ثابت کر سکدے جِنّا چکر ٹال کے۔ ہور وی سانجھاں نیں۔ نارمل !!{proof} اوسے نتیجے دے چکراں والے سبھ توں چھوٹے !!{proof} توں بہت لمّا ہو سکدا اے، جیویں سیکوئنٹ حساب وچ قطعاں والے !!{proof}s قطع توں بناں سبھ توں چھوٹے !!{proof} توں بہت چھوٹے ہو سکدے نیں۔ ذیلی !!{formula} والی خاصیت، کہ تسی ہر نتیجہ ہمیشہ !!{prove} کر سکدے او صرف اوس نتیجے دے ذیلی !!{formula}s نوں اپنے !!{proof} وچ ورت کے، قدرتی اخذ لئی نارمل روپ وچ لیاؤن توں نکلدی اے جیویں سیکوئنٹ حساب لئی قطع ہٹاؤن توں نکلدی اے۔

\end{document}
